\documentclass[11pt,a4paper]{article}
\usepackage[utf8]{inputenc}
\usepackage[T1]{fontenc}
\usepackage{amsmath,amsfonts,amssymb}
\usepackage{graphicx}
\usepackage{hyperref}
\usepackage{listings}
\usepackage[dvipsnames]{xcolor}
\usepackage{booktabs}
\usepackage{float}
\usepackage{geometry}
\geometry{margin=1in}

\definecolor{zenblue}{RGB}{41,121,255}
\definecolor{zengreen}{RGB}{52,199,89}
\definecolor{codegray}{RGB}{245,245,245}

\hypersetup{colorlinks=true,linkcolor=zenblue,urlcolor=zenblue,citecolor=zenblue}

\lstset{
    backgroundcolor=\color{codegray},
    basicstyle=\ttfamily\small,
    breaklines=true,
    captionpos=b,
    frame=single,
    numbers=left,
    numberstyle=\tiny\color{gray}
}

\title{
    \vspace{-2cm}
    \Large \textbf{Zen AI Model Family} \\
    \vspace{0.5cm}
    \Huge \textbf{Zen-Musician} \\
    \vspace{0.3cm}
    \large Packaging YuE for Long-Form Lyrics-to-Song Generation in the Zen Stack \\
    \vspace{0.5cm}
    \normalsize Technical Report v2025.01
}

\author{
    Hanzo AI Research Team\thanks{research@hanzo.ai} \and
    Zoo Labs Foundation\thanks{foundation@zoo.ngo}
}

\date{January 2025}

\begin{document}

\maketitle

\begin{abstract}
\textbf{Zen-Musician} is a packaging and integration of \textbf{YuE}, the open
long-form music generation foundation model developed by \textbf{M-A-P
(Multimodal Art Projection)} and HKUST and released under the permissive
\textbf{Apache-2.0} license \cite{yuan2025yue,yuehf}. We make no claim of a
from-scratch model: Zen-Musician redistributes YuE's published weights
unmodified and exposes them through the Zen tooling, runtime, and API surface.
This report documents the \emph{real} architecture so that integrators understand
what they are deploying. YuE is a lyrics-to-song system: given lyrics together
with a short genre/style description, it generates full songs containing both
sung vocals and accompaniment. It is built on a decoder-only LLaMA-2-style
autoregressive transformer that operates over discrete audio tokens. The flagship
Stage-1 checkpoint (e.g. \texttt{m-a-p/YuE-s1-7B-anneal-en-cot}) has
approximately 7B parameters and performs track-decoupled next-token prediction to
produce vocal and accompaniment token streams; a separate Stage-2 model
(approximately 1B parameters) refines and upsamples those tokens toward the final
waveform \cite{yuegh}. According to its authors, YuE can generate up to roughly
five minutes of music while maintaining lyrical alignment, coherent musical
structure, and vocal melodies with appropriate accompaniment, and supports style
transfer and reference conditioning through an in-context-learning (ICL) variant
\cite{yuan2025yue}. This document supersedes an earlier draft that incorrectly
described a from-scratch ``Zen MoDE'' architecture and reported fabricated
benchmarks; those claims were false and have been removed.
\end{abstract}

\tableofcontents
\newpage

\section{Introduction}

\subsection{What Zen-Musician Is (and Is Not)}

Zen-Musician is the Zen stack's integration of \textbf{YuE}
(\href{https://github.com/multimodal-art-projection/YuE}{multimodal-art-projection/YuE}),
an open foundation model for full-song generation released by M-A-P and HKUST
under the \textbf{Apache-2.0} license \cite{yuehf,yuan2025yue}. The Apache-2.0
license permits redistribution, modification, and commercial use subject to its
attribution and notice requirements; YuE's authors explicitly encourage creators
to incorporate model outputs, with attribution to ``YuE by HKUST/M-A-P''
recommended \cite{yuehf}.

To be explicit about provenance:

\begin{itemize}
    \item Zen-Musician \textbf{does not} introduce a new architecture. There is
    no ``Zen MoDE,'' no ``Section-Aware Transformer,'' and no unified
    MIDI/audio model. Those concepts appeared in a prior draft of this report and
    were not real.
    \item Zen-Musician \textbf{does not} report its own benchmark numbers. The FAD,
    MOS, genre-consistency, structure-F1, and accompaniment-latency tables in the
    earlier draft were fabricated and have been deleted. Where this report
    mentions evaluation, it attributes claims to YuE's authors and does not invent
    metrics.
    \item Zen-Musician \textbf{is} a redistribution of YuE's published weights,
    wrapped with Zen tooling (a consistent loader, runtime integration, and the
    Zen API surface) so that YuE can be used alongside other Zen models.
\end{itemize}

\subsection{Correcting the README}

A previous Zen-Musician description listed the model as ``1B'' parameters and as a
MIDI/symbolic system. Both were wrong. The packaged flagship is the
\textbf{7B Stage-1} checkpoint; the ``1B'' figure corresponds only to the
\emph{Stage-2} refinement/upsampling model, not the main generator. YuE generates
\emph{audio} (sung vocals plus accompaniment) from \emph{lyrics and genre text}.
It is not a MIDI model and does not generate symbolic scores.

\subsection{The Problem YuE Addresses}

Open music models have historically struggled with \emph{long-form} generation:
producing a coherent multi-minute song with intelligible, aligned lyrics, a
sensible verse/chorus structure, and a singing voice that stays on top of a
suitable accompaniment. YuE targets this lyrics-to-song setting directly, scaling
an autoregressive transformer over audio tokens to song-length outputs
\cite{yuan2025yue}.

\subsection{Model Overview}

\begin{table}[H]
\centering
\begin{tabular}{ll}
\toprule
\textbf{Property} & \textbf{Value} \\
\midrule
Upstream model & YuE (M-A-P / HKUST) \\
License & Apache-2.0 \cite{yuehf} \\
Task & Lyrics + genre $\rightarrow$ full song (vocals + accompaniment) \\
Backbone & Decoder-only, LLaMA-2-style autoregressive transformer \cite{yuan2025yue,touvron2023llama2} \\
Stage-1 generator & $\approx$7B parameters (e.g. \texttt{YuE-s1-7B-anneal-en-cot}) \\
Stage-2 refiner & $\approx$1B parameters (token refinement / upsampling) \\
Representation & Discrete audio tokens (xcodec-style tokenizer) \\
Languages & English, Mandarin, Cantonese, Japanese, Korean (checkpoint-specific) \\
Conditioning & Lyrics with section tags; genre/style tags; optional audio reference (ICL) \\
Reported length & Up to $\approx$5 minutes of music \cite{yuan2025yue} \\
\bottomrule
\end{tabular}
\caption{Zen-Musician = packaged YuE: upstream specifications}
\end{table}

\section{Architecture (as published by YuE's authors)}

The description below summarizes YuE's design from its paper and model
documentation \cite{yuan2025yue,yuehf,yuegh}. Zen-Musician does not alter this
architecture.

\subsection{Audio Tokenization}

YuE operates over \emph{discrete audio tokens} rather than raw waveforms. Audio is
encoded into token sequences by an xcodec-style neural audio codec; the language
model predicts tokens, and the decoder side reconstructs audio from them
\cite{yuehf}. This lets a standard autoregressive transformer model music the same
way a text LLM models language.

\subsection{Stage-1: LLaMA-2-style Autoregressive Generator}

The Stage-1 model is a \textbf{decoder-only, LLaMA-2-style transformer}
\cite{yuan2025yue,touvron2023llama2} with approximately \textbf{7B parameters}.
Its central design choice, as described by the authors, is \textbf{track-decoupled
(``dual-track'') next-token prediction}: rather than modeling a single dense
mixture signal, it predicts \emph{vocal} and \emph{accompaniment} token streams in
a decoupled fashion, which the authors report helps the model keep the singing
voice coherent over long contexts \cite{yuan2025yue}. The authors further describe
\emph{structural progressive conditioning} to maintain lyric alignment across
long-form generation \cite{yuan2025yue}.

\subsection{Stage-2: Refinement / Upsampling}

A separate \textbf{Stage-2} model (approximately \textbf{1B parameters}) takes the
Stage-1 token output and performs residual refinement / upsampling toward the
final audio, improving fidelity before waveform reconstruction \cite{yuegh,yuehf}.
This is the source of the ``1B'' figure that an earlier Zen description mistook for
the size of the whole system.

\subsection{Conditioning Interface}

YuE is prompted with three kinds of input \cite{yuehf}:

\begin{itemize}
    \item \textbf{Genre / style tags}: a short free-text description (the project
    suggests including genre, instrumentation, mood, vocal gender, and timbre).
    \item \textbf{Lyrics}: the song's lyrics, annotated with structural labels such
    as \texttt{[verse]} and \texttt{[chorus]}, organized into sessions/segments.
    \item \textbf{Optional audio reference} (ICL variants): a reference clip used
    for style transfer or voice-style conditioning via in-context learning
    \cite{yuan2025yue}.
\end{itemize}

\subsection{Checkpoint Family}

YuE ships multiple Stage-1 checkpoints, differentiated by language coverage and by
mode --- a chain-of-thought (\texttt{cot}) variant and an in-context-learning
(\texttt{icl}) variant for reference-conditioned generation \cite{yuehf}. Example
flagship checkpoint: \texttt{m-a-p/YuE-s1-7B-anneal-en-cot}.

\section{Capabilities Reported Upstream}

The following are claims made by YuE's authors \cite{yuan2025yue}; Zen-Musician
inherits them by redistributing the same weights and does not independently
re-benchmark them in this report.

\begin{itemize}
    \item \textbf{Long-form songs}: generation of up to roughly five minutes of
    music with maintained lyrical alignment and coherent musical structure.
    \item \textbf{Vocals with accompaniment}: engaging vocal melodies sung over an
    appropriate instrumental backing, produced by the dual-track scheme.
    \item \textbf{Style transfer / reference conditioning}: an in-context-learning
    mode enabling style transfer (e.g. genre conversion while preserving
    accompaniment) and reference-based generation.
    \item \textbf{Multilingual coverage}: checkpoints for English, Mandarin,
    Cantonese, Japanese, and Korean.
\end{itemize}

The YuE paper reports that the system ``matches or even surpasses some of the
proprietary systems in musicality and vocal agility,'' and evaluates music
\emph{understanding} on the MARBLE benchmark \cite{yuan2025yue}. We reproduce this
as an attributed upstream claim only; no numeric scores are asserted here, and
none should be attributed to Zen.

\section{Integration in the Zen Stack}

Zen-Musician's contribution is purely \textbf{packaging and integration}, not
modeling:

\begin{itemize}
    \item \textbf{Distribution}: the YuE Stage-1 and Stage-2 weights are
    redistributed under Apache-2.0 with upstream attribution and license/notice
    files preserved.
    \item \textbf{Loader and runtime}: a consistent entry point so YuE can be
    invoked the same way as other Zen models.
    \item \textbf{API surface}: a thin wrapper exposing YuE's lyrics+genre
    conditioning and two-stage pipeline.
\end{itemize}

\subsection{Resource Notes}

YuE is a 7B-class generator over long audio-token contexts; full-song generation
is GPU-memory intensive, and the upstream project documents substantial VRAM
requirements for long outputs as well as multi-stage (Stage-1 then Stage-2)
inference \cite{yuehf}. Integrators should budget accordingly rather than expecting
real-time, low-memory operation.

\section{Usage}

The snippet below is illustrative of the lyrics-to-song interface. Refer to the
upstream YuE repository \cite{yuegh} for the authoritative inference pipeline,
prompt format, and Stage-1/Stage-2 invocation.

\begin{lstlisting}[language=Python, caption=Zen-Musician (packaged YuE) lyrics-to-song generation]
from zen import ZenMusician

# Loads the packaged YuE Stage-1 (7B) generator; Stage-2 refinement
# runs as part of the pipeline.
model = ZenMusician.from_pretrained("zenlm/zen-musician")  # wraps m-a-p/YuE-s1-7B-*

genre = "inspiring pop female vocal bright piano uplifting"

lyrics = """
[verse]
Lines of the verse go here
Another line that scans and rhymes
[chorus]
The hook that repeats and stays
Carrying the song through its days
"""

song = model.generate(
    genre_tags=genre,
    lyrics=lyrics,
    # optional reference audio enables ICL style/voice conditioning
    # reference_audio="reference.wav",
)
song.save("song.wav")
\end{lstlisting}

\section{Provenance and Attribution}

Zen-Musician redistributes YuE. Required and recommended attribution:

\begin{itemize}
    \item \textbf{Model}: YuE --- Open Full-Song Music Generation Foundation Model.
    \item \textbf{Authors / affiliations}: M-A-P (Multimodal Art Projection) and
    HKUST (R.\ Yuan, H.\ Lin, S.\ Guo, G.\ Zhang, et al.; Y.\ Guo)
    \cite{yuan2025yue}.
    \item \textbf{License}: Apache-2.0 \cite{yuehf}.
    \item \textbf{Suggested credit}: ``YuE by HKUST/M-A-P.''
    \item \textbf{Source}: \href{https://github.com/multimodal-art-projection/YuE}{github.com/multimodal-art-projection/YuE};
    weights at \href{https://huggingface.co/m-a-p}{huggingface.co/m-a-p}.
\end{itemize}

\section{Conclusion}

Zen-Musician is the Zen stack's packaging of YuE, an Apache-2.0 long-form
lyrics-to-song model from M-A-P and HKUST. The real system is a decoder-only,
LLaMA-2-style autoregressive transformer (Stage-1 $\approx$7B) that performs
track-decoupled next-token prediction over discrete audio tokens to generate
vocals and accompaniment, followed by a $\approx$1B Stage-2 refinement/upsampling
model. Conditioning is by lyrics and genre/style text, with an optional ICL mode
for reference-based style transfer. All credit for the architecture, training, and
reported capabilities belongs to YuE's authors; Zen's contribution is integration
and distribution under the upstream license, with attribution. The from-scratch
architecture and benchmark claims in the earlier draft were fabricated and have
been removed.

\begin{thebibliography}{10}
\bibitem{yuan2025yue} R.\ Yuan, H.\ Lin, S.\ Guo, G.\ Zhang, et al., ``YuE: Scaling Open Foundation Models for Long-Form Music Generation,'' arXiv:2503.08638, 2025. \url{https://arxiv.org/abs/2503.08638}
\bibitem{yuehf} M-A-P, ``YuE model cards (Apache-2.0),'' Hugging Face. \url{https://huggingface.co/m-a-p/YuE-s1-7B-anneal-en-cot}
\bibitem{yuegh} Multimodal Art Projection, ``YuE: Open Full-song Music Generation Foundation Model,'' GitHub. \url{https://github.com/multimodal-art-projection/YuE}
\bibitem{touvron2023llama2} H.\ Touvron et al., ``Llama 2: Open Foundation and Fine-Tuned Chat Models,'' arXiv:2307.09288, 2023.
\end{thebibliography}

\end{document}
